% GENERATED FILE. Edit canonical lessons, then run npm run generate:books.
% canonical-source-hash: fnv1a64:c5047c0dd252e6aa
% canonical-lessons: JA-C232-hanashiau, JA-C232-matomeru, JA-C232-makaseru, JA-C232-kotowaru, JA-C232-shiraseru

\chapter{Talking Things Over}
\label{ch:ja-talking-things-over}
\hlchaptermodality{232}

\begin{chapteropening}
\textbf{By the end of this chapter:} \emph{I can talk about discussing, sharing out work and saying no.}

Complete the last lesson of chapter 232: I can talk about discussing, sharing out work and saying no.
\end{chapteropening}

\section[\texorpdfstring{\ja{はなしあう} (hanashiau)}{はなしあう (hanashiau)}]{\ja{はなしあう} (hanashiau) — to discuss; to talk together}

\label{lesson:JA-C232-hanashiau}



\begin{quote}
Before the new one: say the Japanese for a response, then the Japanese for to pass on.
\end{quote}

\subsection*{You'll want to know: \ja{はなしあう}}
\textbf{\ja{はなしあう}} — \emph{hanashiau} — \textquotedblleft{}to discuss; to talk together\textquotedblright{}.

\begin{quote}
\textbf{\ja{どうりょうと} \ja{かいぎで} \ja{はなしあいます。}} — \emph{dōryō to kaigi de hanashiaimasu} — \textquotedblleft{}I discuss things with my colleagues at the meeting.\textquotedblright{}
\end{quote}

Polite \textbf{\ja{はなしあいます}}, \emph{hanashiaimasu}; a godan verb.

\begin{cousinweb}
\textbf{\ja{はなしあう}} (\emph{hanashiau}) is a native Japanese compound (wago), inherited rather than borrowed: it is built from \emph{hanashi}, from \emph{hanasu}, \textquotedblleft{}to speak\textquotedblright{}, and \emph{au}, \textquotedblleft{}to do together\textquotedblright{} (JMdict, KANJIDIC2).
\end{cousinweb}

\begin{grammarlens}[title={What you've built}]
One more word for everyday life at A2.
\end{grammarlens}

\subsection*{Your turn}
\begin{itemize}
\raggedright
  \item \emph{Say it:} \emph{hanashiau}
  \item \emph{Say it:} \emph{hanashiau}, once more
  \item \emph{Say it:} \emph{dōryō to kaigi de hanashiaimasu} — I discuss things with my colleagues at the meeting.
  \item \emph{Recall:} say the Japanese for a response, then the Japanese for to pass on, then say \emph{hanashiau} again
\end{itemize}

\begin{tcolorbox}[breakable,colback=teal!4,colframe=teal!35!black,title={Before you move on}]
What does \ja{はなしあう} mean? (\textquotedblleft{}To discuss; to talk together\textquotedblright{}.) Say it once more.
\end{tcolorbox}
\section[\texorpdfstring{\ja{まとめる} (matomeru)}{まとめる (matomeru)}]{\ja{まとめる} (matomeru) — to put together; to sum up}

\label{lesson:JA-C232-matomeru}



\begin{quote}
Before the new one: say the Japanese for to pass on, then the Japanese for to discuss.
\end{quote}

\subsection*{You'll want to know: \ja{まとめる}}
\textbf{\ja{まとめる}} — \emph{matomeru} — \textquotedblleft{}to put together; to sum up\textquotedblright{}.

\begin{quote}
\textbf{\ja{せんぱいは} \ja{かいぎで} \ja{みんなの} \ja{しつもんを} \ja{まとめます。}} — \emph{senpai wa kaigi de minna no shitsumon o matomemasu} — \textquotedblleft{}My senior sums up everyone's questions at the meeting.\textquotedblright{}
\end{quote}

Polite \textbf{\ja{まとめます}}, \emph{matomemasu}; an ichidan verb.

\begin{cousinweb}
\textbf{\ja{まとめる}} (\emph{matomeru}) is a native Japanese word (wago), inherited rather than borrowed from Chinese; its kanji was attached later as a kun'yomi, and the word is usually written in kana (JMdict, KANJIDIC2).
\end{cousinweb}

\begin{grammarlens}[title={What you've built}]
One more word for everyday life at A2.
\end{grammarlens}

\subsection*{Your turn}
\begin{itemize}
\raggedright
  \item \emph{Say it:} \emph{matomeru}
  \item \emph{Say it:} \emph{matomeru}, once more
  \item \emph{Say it:} \emph{senpai wa kaigi de minna no shitsumon o matomemasu} — My senior sums up everyone's questions at the meeting.
  \item \emph{Recall:} say the Japanese for to pass on, then the Japanese for to discuss, then say \emph{matomeru} again
\end{itemize}

\begin{tcolorbox}[breakable,colback=teal!4,colframe=teal!35!black,title={Before you move on}]
What does \ja{まとめる} mean? (\textquotedblleft{}To put together; to sum up\textquotedblright{}.) Say it once more.
\end{tcolorbox}
\section[\texorpdfstring{\ja{まかせる} (makaseru)}{まかせる (makaseru)}]{\ja{まかせる} (makaseru) — to entrust; to leave to}

\label{lesson:JA-C232-makaseru}



\begin{quote}
Before the new one: say the Japanese for to discuss, then the Japanese for to put together.
\end{quote}

\subsection*{You'll want to know: \ja{まかせる}}
\textbf{\ja{まかせる}} — \emph{makaseru} — \textquotedblleft{}to entrust; to leave to\textquotedblright{}.

\begin{quote}
\textbf{\ja{しゃちょうは} \ja{どうりょうに} \ja{しごとを} \ja{まかせます。}} — \emph{shachō wa dōryō ni shigoto o makasemasu} — \textquotedblleft{}The company president leaves the work to my colleague.\textquotedblright{}
\end{quote}

Polite \textbf{\ja{まかせます}}, \emph{makasemasu}; an ichidan verb.

\begin{cousinweb}
\textbf{\ja{まかせる}} (\emph{makaseru}) is a native Japanese word (wago), inherited rather than borrowed from Chinese; its kanji was attached later as a kun'yomi (JMdict, KANJIDIC2).
\end{cousinweb}

\begin{grammarlens}[title={What you've built}]
One more word for everyday life at A2.
\end{grammarlens}

\subsection*{Your turn}
\begin{itemize}
\raggedright
  \item \emph{Say it:} \emph{makaseru}
  \item \emph{Say it:} \emph{makaseru}, once more
  \item \emph{Say it:} \emph{shachō wa dōryō ni shigoto o makasemasu} — The company president leaves the work to my colleague.
  \item \emph{Recall:} say the Japanese for to discuss, then the Japanese for to put together, then say \emph{makaseru} again
\end{itemize}

\begin{tcolorbox}[breakable,colback=teal!4,colframe=teal!35!black,title={Before you move on}]
What does \ja{まかせる} mean? (\textquotedblleft{}To entrust; to leave to\textquotedblright{}.) Say it once more.
\end{tcolorbox}
\section[\texorpdfstring{\ja{ことわる} (kotowaru)}{ことわる (kotowaru)}]{\ja{ことわる} (kotowaru) — to refuse; to turn down}

\label{lesson:JA-C232-kotowaru}



\begin{quote}
Before the new one: say the Japanese for to put together, then the Japanese for to entrust.
\end{quote}

\subsection*{You'll want to know: \ja{ことわる}}
\textbf{\ja{ことわる}} — \emph{kotowaru} — \textquotedblleft{}to refuse; to turn down\textquotedblright{}.

\begin{quote}
\textbf{\ja{あには} \ja{よく} \ja{しごとを} \ja{ことわります。}} — \emph{ani wa yoku shigoto o kotowarimasu} — \textquotedblleft{}My elder brother often turns down work.\textquotedblright{}
\end{quote}

Polite \textbf{\ja{ことわります}}, \emph{kotowarimasu}; a godan verb.

\begin{cousinweb}
\textbf{\ja{ことわる}} (\emph{kotowaru}) is a native Japanese word (wago), inherited rather than borrowed from Chinese; its kanji was attached later as a kun'yomi (JMdict, KANJIDIC2).
\end{cousinweb}

\begin{grammarlens}[title={What you've built}]
One more word for everyday life at A2.
\end{grammarlens}

\subsection*{Your turn}
\begin{itemize}
\raggedright
  \item \emph{Say it:} \emph{kotowaru}
  \item \emph{Say it:} \emph{kotowaru}, once more
  \item \emph{Say it:} \emph{ani wa yoku shigoto o kotowarimasu} — My elder brother often turns down work.
  \item \emph{Recall:} say the Japanese for to put together, then the Japanese for to entrust, then say \emph{kotowaru} again
\end{itemize}

\begin{tcolorbox}[breakable,colback=teal!4,colframe=teal!35!black,title={Before you move on}]
What does \ja{ことわる} mean? (\textquotedblleft{}To refuse; to turn down\textquotedblright{}.) Say it once more.
\end{tcolorbox}
\section[\texorpdfstring{\ja{しらせる} (shiraseru)}{しらせる (shiraseru)}]{\ja{しらせる} (shiraseru) — to inform; to let know}

\label{lesson:JA-C232-shiraseru}



\begin{quote}
Before the new one: say the Japanese for to entrust, then the Japanese for to refuse.
\end{quote}

\subsection*{You'll want to know: \ja{しらせる}}
\textbf{\ja{しらせる}} — \emph{shiraseru} — \textquotedblleft{}to inform; to let know\textquotedblright{}.

\begin{quote}
\textbf{\ja{せんせいは} \ja{がくせいに} \ja{しけんの} \ja{じかんを} \ja{しらせます。}} — \emph{sensei wa gakusei ni shiken no jikan o shirasemasu} — \textquotedblleft{}The teacher tells the students the time of the exam.\textquotedblright{}
\end{quote}

Polite \textbf{\ja{しらせます}}, \emph{shirasemasu}; an ichidan verb.

\begin{cousinweb}
\textbf{\ja{しらせる}} (\emph{shiraseru}) is a native Japanese word (wago), inherited rather than borrowed: it is built on \emph{shiru}, \textquotedblleft{}to know\textquotedblright{}, and means \textquotedblleft{}to make someone know\textquotedblright{} (JMdict, KANJIDIC2).
\end{cousinweb}

\begin{grammarlens}[title={What you've built}]
One more word for everyday life at A2.
\end{grammarlens}

\subsection*{Your turn}
\begin{itemize}
\raggedright
  \item \emph{Say it:} \emph{shiraseru}
  \item \emph{Say it:} \emph{shiraseru}, once more
  \item \emph{Say it:} \emph{sensei wa gakusei ni shiken no jikan o shirasemasu} — The teacher tells the students the time of the exam.
  \item \emph{Recall:} say the Japanese for to entrust, then the Japanese for to refuse, then say \emph{shiraseru} again
\end{itemize}

\begin{tcolorbox}[breakable,colback=teal!4,colframe=teal!35!black,title={Before you move on}]
What does \ja{しらせる} mean? (\textquotedblleft{}To inform; to let know\textquotedblright{}.) Say it once more.
\end{tcolorbox}
